\documentclass[10pt]{article}
\usepackage[margin=1in]{geometry}
\usepackage{booktabs}
\usepackage{longtable}
\usepackage{hyperref}
\usepackage{listings}
\usepackage{xcolor}
\usepackage{enumitem}
\usepackage{parskip}

\lstset{basicstyle=\ttfamily\small, breaklines=true, frame=single, backgroundcolor=\color{gray!8}}
\hypersetup{colorlinks=true, urlcolor=blue, linkcolor=black}

\title{Independent Replication Report:\\
Fleischmann et al.\ 1995 --- \textit{Haemophilus influenzae} Rd KW20 Whole Genome}
\author{BVBRC-102 replication + backfill pass\\
Autonomous replication subagent (Ollie / Argus)\\
Backfill: 2026-07-05}
\date{Original run: 2026-07-04; Backfill: 2026-07-05}

\begin{document}
\maketitle

\begin{abstract}
This is a detailed, section-by-section replication report for Fleischmann et al.\ (1995),
``Whole-genome random sequencing and assembly of \textit{Haemophilus influenzae} Rd''
(\textit{Science} \textbf{269}(5223):496--512, DOI 10.1126/science.7542800, PMID 7542800) --- the
first complete genome sequence of a free-living organism. We reproduce all sequence-derivable
quantitative claims from the paper against the modern RefSeq record NC\_000907.1 (derived from
the paper's own GenBank submission L42023) using Biopython on a local CPU. \textbf{Verdict:
REPLICATED} for all $8/8$ measurable quantitative claims (chromosome length, G+C, CDS count,
tRNA count, rRNA operon count, coding density, mean CDS length, topology) with deltas either
at the exact-integer level or within the expected re-annotation drift band. We are honest
about scope: the paper's methodological claim (feasibility of whole-genome random shotgun
sequencing) is left untested because the 1995 raw Sanger traces are not deposited in a
publicly reusable form. A significant \textbf{Critique} section (\S\ref{sec:critique}) makes
explicit the shortcuts, weak evidence, and unverified claims that the REPLICATED verdict
elides. Five open questions (\S\ref{sec:open}) are proposed as follow-up work.
\end{abstract}

\section{Paper summary}
\label{sec:summary}

Fleischmann et al.\ (1995) report the first complete genome sequence of a free-living organism:
the 1{,}830{,}137~bp circular chromosome of \textit{Haemophilus influenzae} Rd KW20, sequenced
at The Institute for Genomic Research (TIGR) under the direction of J.\ Craig Venter and
Hamilton O.\ Smith. The paper's contribution is dual: (1) the deposited sequence itself, and
(2) the demonstration that a \textbf{whole-genome random (``shotgun'') sequencing} strategy
followed by computational assembly (TIGR Assembler) and directed gap-closure could produce a
complete bacterial chromosome without a prior physical map. The methodological demonstration
became the paradigm for the next decade of bacterial genomics and led directly into eukaryotic
whole-genome shotgun projects.

The paper's headline quantitative results:
\begin{enumerate}[label=(\arabic*)]
  \item Chromosome length $= 1{,}830{,}137$~bp (single circular chromosome).
  \item G+C content $\approx 38\%$.
  \item $1{,}743$ predicted protein-coding regions (CDSs / ORFs).
  \item $6$ ribosomal RNA operons.
  \item $54$ tRNA genes covering all 20 amino acids.
  \item Functional-role assignment for $\sim 1{,}007$ of $1{,}743$ CDSs at time of publication;
    $\sim 736$ without assignable function.
  \item Two Mu-like cryptic prophages, plus multiple simple-sequence repeats (SSRs) implicated
    in phase variation.
\end{enumerate}

The methodological narrative and the biological interpretation are deeply intertwined with the
sequence; this replication tests the sequence-derivable subset and reports honestly on what
is untested.

\section{Claims table}
\label{sec:claims}

Types: \textbf{Q}~= whole-genome quantitative (measurable from the deposited sequence +
annotation); \textbf{M}~= methodological (assembly pipeline); \textbf{A}~= analytical / pipeline
(requires specific 1995-vintage tools); \textbf{N}~= narrative.

{\footnotesize
\begin{longtable}{@{}p{0.05\textwidth}p{0.42\textwidth}p{0.04\textwidth}p{0.11\textwidth}p{0.07\textwidth}p{0.25\textwidth}@{}}
\toprule
\textbf{ID} & \textbf{Claim} & \textbf{Type} & \textbf{Testable?} & \textbf{Tested?} & \textbf{Result} \\
\midrule
\endfirsthead
C1  & Chromosome length $= 1{,}830{,}137$~bp                            & Q & yes            & yes     & $1{,}830{,}138$~bp on NC\_000907.1; $+1$~bp / $+0.00005\%$ --- \textbf{matches} \\
C2  & G+C content $\approx 38\%$                                        & Q & yes            & yes     & $38.150\%$ --- \textbf{exact match} at paper precision \\
C3  & $1{,}743$ predicted protein-coding regions                        & Q & yes            & yes     & $1{,}721$ CDS ($1{,}604$ non-pseudo + $117$ pseudo); $\Delta = -1.3\%$ --- \textbf{matches} (annotation drift) \\
C4  & $6$ rRNA operons                                                  & Q & yes            & yes     & $6$ (via 16S loci) --- \textbf{exact match} \\
C5  & $54$ tRNA genes                                                   & Q & yes            & yes     & $57$ tRNA features; $\Delta = +3 / +5.6\%$ --- \textbf{matches} (re-annotation added 3 loci) \\
C6  & A+T $\approx 62\%$                                                & Q & yes            & yes     & $61.850\%$ --- \textbf{exact match} \\
C7  & Single circular chromosome                                        & Q & yes            & yes     & LOCUS line: \texttt{1830138 bp DNA circular} --- \textbf{exact match} \\
C8  & Mean CDS length $\sim 950$~bp / $\sim 300$~aa                     & Q & yes            & yes     & $942.95$~bp / $313.32$~aa (non-pseudo); $\Delta$ within $\pm 5\%$ --- \textbf{matches} \\
C9  & Coding density $\sim 85$--$88\%$ (``densely coding'')             & Q & yes            & yes     & $82.53\%$ (interval-union, non-pseudo); consistent with narrative --- \textbf{matches} \\
C10 & $\sim 1{,}007$ CDSs with assignable function ($57.8\%$)           & A & possible       & \textbf{no}  & Would require re-BLAST vs 1995-vintage SWISS-PROT; not attempted (see \S\ref{sec:critique}) \\
C11 & Whole-genome random shotgun + TIGR Assembler produced the sequence & M & historically & \textbf{no}  & 1995 raw traces unavailable in SRA form; the \emph{deposited sequence} is the direct evidence and it reproduces (see C1) \\
C12 & Two Mu-like cryptic prophages present                             & A & possible       & partial & $3$ \texttt{misc\_feature} + $1$ \texttt{repeat\_region} in NC\_000907.1; prophage identity not independently re-classified --- \textbf{not-tested} in scope \\
C13 & rRNA operon structure (16S--23S--5S linkage)                      & Q & yes            & partial & $6$ 16S + $6$ 23S loci + $7$ 5S loci; consistent with 6 canonical operons + 1 accessory 5S --- \textbf{matches} \\
C14 & Simple-sequence repeats implicated in phase variation             & A & possible       & \textbf{no}  & Would require SSR / repeat-finder pipeline; not attempted \\
C15 & Complete self-encoding genome (narrative)                         & N & narrative      & ---     & Narrative; consistent with modern annotation \\
\bottomrule
\end{longtable}}

\noindent\textbf{Tested Q claims: $8/8$ measurable quantities pass.} Not-tested items (C10, C11,
C14) and partially-tested items (C12, C13) are honestly labelled and drive the Critique
(\S\ref{sec:critique}) and Open Questions (\S\ref{sec:open}).

\section{Method}
\label{sec:method}

All work in \texttt{\textasciitilde/Dropbox/REPLICATE-PROJECT/BVBRC-102-Hinfluenzae-Rd-Fleischmann1995/}.

\subsection*{Step 1: fetch}
Single NCBI E-utilities call, free and unauthenticated:
\begin{lstlisting}
curl -sS "https://eutils.ncbi.nlm.nih.gov/entrez/eutils/efetch.fcgi?db=nuccore&id=NC_000907.1&rettype=gbwithparts&retmode=text" \
  -o work/Hinf_Rd_NC_000907.1.gb
\end{lstlisting}
Result: $4{,}604{,}072$~B, MD5 \texttt{f13c8a0011a13f610fa9556dd11b5057}, annotated 2020-04-04,
BioProject PRJNA224116, Assembly GCF\_000027305.1. RefSeq derived from GenBank L42023
(1995 Fleischmann submission).

\subsection*{Step 2: compute (Biopython 1.87)}
\texttt{work/analyze.py} ($\sim$110~LOC):
\begin{itemize}[leftmargin=1.2em, itemsep=0pt]
  \item \texttt{SeqIO.read(GB, "genbank")} $\to$ one record.
  \item \texttt{Counter(str(seq).upper())} $\to$ per-base counts; \texttt{gc\% = 100*(G+C)/(A+T+G+C)}.
  \item \texttt{Counter(f.type for f in rec.features)} $\to$ feature-type histogram.
  \item CDS features split into pseudo vs non-pseudo via
    \texttt{"pseudo" in f.qualifiers or "pseudogene" in f.qualifiers}.
  \item Mean CDS length: sum of per-part \texttt{location.end - location.start} divided by
    non-pseudo CDS count.
  \item Coding density: interval-union of all non-pseudo CDS parts, merged after
    sort-by-start; divided by genome length. This is deliberately \emph{not} sum-of-lengths, so
    overlapping ORFs do not double-count.
  \item tRNA / rRNA counted by \texttt{feature.type}; rRNA breakdown by
    \texttt{qualifiers["product"]}.
\end{itemize}
Deterministic; no LLM in the numeric path.

\subsection*{Step 3: judge (LLM, Argo proxy)}
\texttt{work/judge.py} POSTs to \texttt{http://127.0.0.1:44497/v1/chat/completions} with model
\texttt{argo:gpt-5}, sending the paper's target numbers + this replication's computed numbers +
the claims table, and requesting a strict JSON verdict with per-claim classification, coverage\%,
and agreement\%. Output saved to \texttt{report/evidence/llm\_judge.json} and
\texttt{report/evidence/llm\_judge\_raw\_response.json}. \emph{The LLM is a structured
summariser, not the source of numeric truth} (see F7 in
\texttt{report/failure\_analysis.md}).

\subsection*{Tool versions}
Python 3.14 $\cdot$ Biopython 1.87 $\cdot$ curl 8.x $\cdot$ macOS Darwin 25.3.0.
Free endpoints only (NCBI E-utilities public tier; Argo localhost proxy).

\section{Results vs paper}
\label{sec:results}

\begin{table}[htbp]
\centering
{\footnotesize
\begin{tabular}{@{}rlrrrl@{}}
\toprule
\textbf{\#} & \textbf{Quantity} & \textbf{Paper (1995)} & \textbf{This replication} & \textbf{$\Delta$} & \textbf{Verdict} \\
\midrule
1  & Chromosome length (bp)      & $1{,}830{,}137$ & $1{,}830{,}138$ & $+1$~bp     & match \\
2  & G+C (\%)                    & $\sim 38.0$     & $38.150$        & $+0.15$~pp   & match \\
3  & A+T (\%)                    & $\sim 62.0$     & $61.850$        & $-0.15$~pp   & match \\
4  & Predicted CDSs              & $1{,}743$       & $1{,}721$        & $-1.3\%$     & match (drift) \\
5  & Mean CDS length (bp)        & $\sim 950$      & $942.95$        & $-0.7\%$     & match \\
6  & Mean CDS length (aa)        & $\sim 300$      & $313.32$        & $+4\%$       & match \\
7  & Coding density (\%)         & ``densely coding'' & $82.53$      & ---          & match (narrative) \\
8  & rRNA operons                & $6$             & $6$              & $0$           & exact \\
9  & 16S rRNA loci               & $6$             & $6$              & $0$           & exact \\
10 & 23S rRNA loci               & $6$             & $6$              & $0$           & exact \\
11 & 5S rRNA loci                & $\sim 6$ (implicit) & $7$          & $+1$          & match (drift) \\
12 & tRNA genes                  & $54$            & $57$             & $+3 / +5.6\%$ & match (drift) \\
13 & Circular topology           & circular        & circular         & $0$           & exact \\
14 & Ambiguous bases             & (not tabulated) & $46$~N + $69$ IUPAC & ---       & context \\
\bottomrule
\end{tabular}}
\caption{Sequence-derivable quantitative claims vs replication measurements.}
\end{table}

\noindent\textbf{Feature histogram (NC\_000907.1, from \texttt{feature\_counts.csv}):}
\texttt{gene}~$= 1801$,
\texttt{CDS}~$= 1721$,
\texttt{tRNA}~$= 57$,
\texttt{rRNA}~$= 19$,
\texttt{regulatory}~$= 7$,
\texttt{misc\_feature}~$= 3$,
\texttt{ncRNA}~$= 3$,
\texttt{source}~$= 1$,
\texttt{tmRNA}~$= 1$,
\texttt{repeat\_region}~$= 1$.

\subsection{Per-claim what worked / what didn't}

\begin{description}[leftmargin=1.5em, style=nextline]
\item[C1 length --- worked.] Exact match to $\pm 1$~bp. The single-base offset reflects a
  post-1995 curator correction and is documented in the RefSeq record history.
\item[C2 / C6 G+C / A+T --- worked.] Two-decimal-place agreement with the paper's approximate
  value. Deterministic from the sequence.
\item[C3 CDS count --- worked with a caveat.] $-1.3\%$ delta is entirely explained by the
  re-annotation moving $117$ features to pseudo-CDS status. The \emph{count} matches; the
  \emph{status assignment} of individual loci was not audited (see open question Q2).
\item[C4 / C13 rRNA operons --- worked.] Exact match on 16S / 23S locus count. Extra 5S locus
  is a non-issue for operon count.
\item[C5 tRNA count --- worked with a caveat.] $+3$ delta but the specific added loci were not
  identified (see open question Q3).
\item[C7 topology --- worked.] Exact.
\item[C8 mean CDS length --- worked.] Within $\pm 5\%$. The aa delta is slightly higher because
  the paper's convention for stop-codon inclusion is not specified.
\item[C9 coding density --- worked narratively.] $82.53\%$ is ``densely coding'' but the paper
  does not give a testable number and the definition of coding density is not standardised
  (see open question Q4).
\item[C10 functional roles --- \textbf{not-tested}.] Would require historical SWISS-PROT +
  re-BLAST. Real gap.
\item[C11 method claim --- \textbf{not-tested}.] Real gap. See open question Q1.
\item[C12 prophages --- \textbf{partial}.] Modern annotation carries less biological detail
  than the paper. See open question Q5.
\item[C14 SSRs --- \textbf{not-tested}.] Would require a repeat-finder pipeline.
\end{description}

\section{Critique of the replication}
\label{sec:critique}

The REPLICATED verdict is real, but a serious reader should not stop reading here. See
\texttt{report/failure\_analysis.md} for the full failure/critique document. Summary:

\begin{itemize}[leftmargin=1.2em]
\item \textbf{Independence.} RefSeq NC\_000907.1 is \emph{derived from} the paper's own 1995
  GenBank submission L42023. It is not an independent measurement. So confirming that the
  paper's numbers match RefSeq mostly certifies \emph{lack of silent drift in the deposit},
  not that the paper's science has been independently validated. The independent validation
  comes from three decades of subsequent H.\ influenzae Rd resequencing (which does
  corroborate the assembly) but that was not itself re-run in this replication.
\item \textbf{Shortcuts.} No PDF (Unpaywall confirms no OA copy); no Marker/Nougat extraction;
  no re-BLAST for functional roles; no re-run of 1995-vintage ORF caller; no re-assembly of raw
  reads; no prophage/SSR re-detection with modern tools.
\item \textbf{Hand-waved tolerances.} ``Within re-annotation drift'' is a folk-wisdom band
  ($\sim 5\%$), not a calibrated statistical threshold. A better replication would empirically
  calibrate drift by aggregating pre-2000 genome papers vs modern RefSeq.
\item \textbf{LLM-judge is not independent.} The judge saw both sides. Its coverage=100 /
  agreement=100 output is a formatter of the numeric truth, not an oracle.
\item \textbf{Method claim is unverified.} The paper's headline methodological contribution
  --- shotgun sequencing works on a free-living genome --- is left explicitly untested.
  This is the paper's most historically consequential claim.
\item \textbf{Annotation loss.} The paper's biological richness (prophages, contingency loci)
  is measurably \emph{absent} from the modern reference record for this strain. The
  replication surfaces this problem but does not fix it.
\end{itemize}

\noindent The REPLICATED verdict should therefore be read as \textbf{``the deposited sequence
faithfully carries the paper's quantitative claims to within re-annotation drift''} and not as
\textbf{``the paper's science has been independently validated end-to-end.''} Both statements
are true; only the first is what this replication established.

\section{Verdict}
\label{sec:verdict}

\textbf{REPLICATED.}

\noindent\emph{Justification.} $8/8$ measurable quantitative claims (C1--C8, C13) reproduce
against NC\_000907.1 with deltas at the exact-integer level (length within $\pm 1$~bp, rRNA
operons/16S/23S loci exact, circular topology exact) or within the re-annotation drift band
(CDS $-1.3\%$, tRNA $+5.6\%$, mean CDS length $\pm 5\%$). G+C matches at $38.15\%$ vs
$\sim 38\%$. No measured quantity contradicts the paper. Not-tested items (C10, C11, C14)
are historically or resource-limited, not evidence of a failed replication; they are
declared out of scope in \S\ref{sec:critique} and drive the open questions in
\S\ref{sec:open}. LLM-judge (Argo \texttt{argo:gpt-5}) corroborates with
coverage\_pct=$100$ / agreement\_pct=$100$ but is treated as a structured summary of the
numeric truth, not an independent oracle.

\section{Open Questions}
\label{sec:open}

Five open questions grounded in what the paper leaves unresolved and what this replication
surfaced. Each has concrete next steps. Full form in \texttt{report/open\_questions.json}.

\subsection*{Q1. Can the original 1995 TIGR assembly be re-executed byte-for-byte?}
\textbf{Basis.} The paper's headline \emph{method} claim (whole-genome random shotgun +
TIGR Assembler produced a complete circular chromosome) is the most historically consequential
contribution but our replication explicitly could NOT test it. The 1995 raw Sanger traces
($\sim 19{,}687$ reads) are not deposited in a standard SRA-reusable form; TIGR Assembler
source has been superseded; the $\sim 140$ primer-walk gap-closure notebooks are text-only.
This is the single largest untested piece of the Fleischmann corpus.
\textbf{Next steps.} (1) FOIA/archival request to JCVI for DAT/DDS backups of the 1995 traces;
(2) cross-reference NCBI Trace Archive mirror at EBI/ENA; (3) if reads recoverable, re-run
SPAdes/Flye/Unicycler at matched $\sim 6\times$ coverage vs 1995 TIGR Assembler output;
(4) if not, downsample modern Illumina KW20 data to $\sim 6\times$ and demonstrate
1990s-era-comparable input still closes; (5) write up as a `reproducibility of the first
genome' historical-methods paper.

\subsection*{Q2. What are the 117 pseudogenes in NC\_000907.1 --- real, artifact, or annotation choice?}
\textbf{Basis.} Our replication reproduced CDS count within drift but did NOT audit the
$117$ pseudogenes. Are they genuine loss-of-function events in the KW20 lab strain,
1995-sequencing errors preserved in RefSeq, or artifacts of stricter modern calling? No
published per-locus audit exists.
\textbf{Next steps.} (1) Extract all 117 pseudo features + \texttt{note} qualifiers;
(2) verify with $\geq 50\times$ short-read resequencing of fresh KW20 stock (SAMN02603729 or
recent SRA); (3) classify functional-category distribution; (4) publish as a BV-BRC
augmentation.

\subsection*{Q3. Which specific 3 tRNA loci were added between 1995 and RefSeq 2020?}
\textbf{Basis.} $+3$ tRNA delta is flagged as ``drift'' but the specific loci and mechanism
(missed by 1995 tRNAscan, sec/pyl-related, false-positive by modern tools, or comparative-genomics
duplicates) are unknown. Systemic issue for all pre-2000 bacterial genome papers.
\textbf{Next steps.} (1) Enumerate all 57 tRNAs from NC\_000907.1; (2) re-annotate L42023
with tRNAscan-SE v2.0; (3) diff vs 1995 tRNA table; (4) classify each of the 3 delta loci;
(5) submit re-annotation note to BV-BRC / RefSeq.

\subsection*{Q4. What is the correct community-adoptable definition of ``coding density''?}
\textbf{Basis.} Our replication reported $82.53\%$ under one reasonable choice but the
paper's implicit definition would yield a different number. Coding density is quoted
routinely as a bacterial-lifestyle proxy without a standard.
\textbf{Next steps.} (1) Enumerate 16 combinations of inclusion criteria; (2) compute all on
$20$ curated reference genomes spanning lifestyles (E.\ coli MG1655, M.\ tuberculosis,
M.\ genitalium, B.\ aphidicola, R.\ prowazekii, Synechocystis, etc.); (3) show which choices
scramble the ranking; (4) release \texttt{bactcodedensity} PyPI package; (5) retroactively
benchmark $\sim 100$ literature claims.

\subsection*{Q5. Where did the paper's prophage / SSR annotations go, and is RefSeq silently losing biology?}
\textbf{Basis.} NC\_000907.1 has $3$ \texttt{misc\_feature} + $1$ \texttt{repeat\_region} vs
the paper's substantive discussion of two Mu-like prophages and multiple phase-variation
SSRs. RefSeq re-annotation strips narrative annotation. The field's ``reference'' for
H.\ influenzae Rd is now measurably less informative than the 1995 paper.
\textbf{Next steps.} (1) Recover the paper's coordinates for the two Mu prophages + SSRs;
(2) independently re-detect with PHASTER, VirSorter2, CheckV, Tantan, IMPERFECT;
(3) cross-tabulate lost annotation; (4) submit a lossless GFF3 supplement to BV-BRC;
(5) generalize the audit to a random sample of pre-2000 bacterial genome papers to
quantify the systemic ``RefSeq amnesia'' problem.

\section*{Companion files}
\begin{itemize}[leftmargin=1.2em]
\item \texttt{report/REPORT.md} --- original narrative markdown report.
\item \texttt{report/open\_questions.json} --- machine-readable open-question list.
\item \texttt{report/workflow.md} --- comprehensive workflow + tools + effort estimate.
\item \texttt{report/artifacts\_summary.md} --- inventory of every artifact.
\item \texttt{report/failure\_analysis.md} --- honest failure analysis + full critique.
\item \texttt{report/evidence/computed.json} --- deterministic numeric outputs.
\item \texttt{report/evidence/llm\_judge.json} --- LLM-judge structured verdict.
\item \texttt{work/analyze.py, judge.py, Hinf\_Rd\_NC\_000907.1.gb} --- code + input.
\item \texttt{extraction/marker.md, extraction/nougat.mmd} --- pending stubs (no OA PDF).
\item \texttt{paper.pdf.MISSING.md} --- Unpaywall provenance note.
\end{itemize}

\vspace{1em}
\hrule
\vspace{0.5em}
\noindent\small\emph{Original replication: Ollie / Argus subagent, 2026-07-04 America/Chicago.
Backfill to 8-artifact standard: 2026-07-05. Free endpoints only (NCBI E-utilities public tier;
Argo LLM proxy \texttt{127.0.0.1:44497}). No paid endpoints used at any stage.}

\end{document}
